\documentclass[11pt]{article}

\usepackage[a4paper,margin=1in]{geometry}
\usepackage{amsmath,amssymb}
\usepackage{hyperref}

\setlength{\parindent}{0pt}
\setlength{\parskip}{0.7em}

\title{Inhomogeneous Couplings}
\author{}
\date{}

\begin{document}
\maketitle

This note describes the physics and high-level implementation rules for adding
inhomogeneous couplings to the custom MCWF trajectory code.

The homogeneous implementation uses one active-manifold sector label,

\begin{equation}
N_J=N_\downarrow+N_e,
\end{equation}

and one Dicke-basis label inside each sector,

\begin{equation}
|n_e\rangle,
\qquad
n_e=0,\dots,N_J.
\end{equation}

With inhomogeneous couplings, the atoms are split into coupling groups. Global
permutation symmetry is broken, but permutation symmetry remains inside each
group.

\section*{Split the atoms into two coupling groups}

Split the full atom ensemble into two coupling groups,

\begin{equation}
N_1+N_2=N.
\end{equation}

Here \(N_1\) is the number of physical atoms with coupling \(\omega_1\), and
\(N_2\) is the number of physical atoms with coupling \(\omega_2\).

The high-level user inputs should be

\begin{equation}
N,
\qquad
dN,
\qquad
N_1,
\qquad
N_2,
\qquad
\omega_1.
\end{equation}

Set

\begin{equation}
\omega_2
=
\frac{N-N_1\omega_1}{N_2},
\qquad
N=N_1+N_2.
\end{equation}

Do not require the user to pass \(\omega_2\) separately. It is computed once
from the physical group sizes using the atom-number weighted normalization

\begin{equation}
N_1\omega_1+N_2\omega_2=N.
\end{equation}

This keeps the average coupling fixed when comparing homogeneous and
inhomogeneous runs. When \(N_2=0\), \(\omega_2\) multiplies an empty physical
group and is irrelevant.

The homogeneous limit is

\begin{equation}
\omega_1=\omega_2=1.
\end{equation}

\section*{Group-resolved strong-symmetry sectors}

The strong-symmetry sector is no longer only a scalar \(N_J\). In inhomogeneous
mode, use the pair

\begin{equation}
(N_{J,1},N_{J,2})
\end{equation}

as the sector label.

Here \(N_{J,1}\) is the number of active-manifold atoms in group 1, and
\(N_{J,2}\) is the number of active-manifold atoms in group 2.

The total active-manifold number is

\begin{equation}
N_J=N_{J,1}+N_{J,2}.
\end{equation}

This remains the physical number of atoms in the active manifold. The coupling
weights \(\omega_1\) and \(\omega_2\) affect the drive and collective jump
operators, but they do not redefine the sector label.

The allowed group-resolved sectors obey

\begin{equation}
0\leq N_{J,1}\leq N_1,
\qquad
0\leq N_{J,2}\leq N_2.
\end{equation}

Inside each group-resolved sector, the basis is

\begin{equation}
|n_{e,1},n_{e,2}\rangle,
\end{equation}

where

\begin{equation}
n_{e,1}=0,\dots,N_{J,1},
\qquad
n_{e,2}=0,\dots,N_{J,2}.
\end{equation}

The sector dimension is therefore

\begin{equation}
(N_{J,1}+1)(N_{J,2}+1).
\end{equation}

This replaces the homogeneous sector basis dimension \(N_J+1\).

\section*{High-level sector construction}

Do not require the user to manually pass all low-level group-resolved sector
coefficients.

The low-level representation may remain

\begin{verbatim}
{(Nj1, Nj2): coeff}
\end{verbatim}

but user-facing inhomogeneous runs should be able to pass

\begin{verbatim}
N, dN, N1, N2, omega1
\end{verbatim}

and have the code construct all allowed group-resolved sectors automatically.

First choose the total active-manifold sectors exactly as in the homogeneous
code:

\begin{equation}
N_J\in\{N/2-dN,\dots,N/2+dN\}.
\end{equation}

For each target homogeneous-sector label \(N_J\), include every valid pair

\begin{equation}
(N_{J,1},N_{J,2})
\end{equation}

satisfying

\begin{equation}
N_{J,1}+N_{J,2}=N_J,
\end{equation}

with

\begin{equation}
0\leq N_{J,1}\leq N_1,
\qquad
0\leq N_{J,2}\leq N_2.
\end{equation}

Example:

\begin{equation}
N=10,
\qquad
dN=1,
\qquad
N_1=3,
\qquad
N_2=7.
\end{equation}

The selected total sectors are

\begin{equation}
N_J=4,5,6.
\end{equation}

The generated group-resolved sectors are therefore the valid pairs whose
physical active-manifold count matches one of these values.

For example, the generated group-resolved sectors should be

\begin{align}
N_J=4: &\qquad (0,4),(1,3),(2,2),(3,1),\\
N_J=5: &\qquad (0,5),(1,4),(2,3),(3,2),\\
N_J=6: &\qquad (0,6),(1,5),(2,4),(3,3).
\end{align}

So the helper should return sector coefficients for all 12 group-resolved
sectors.

The default internal state in each group-resolved sector should be

\begin{equation}
|n_{e,1}=0,n_{e,2}=0\rangle,
\end{equation}

meaning all active atoms in both groups start in \(|\downarrow\rangle\).

\section*{Initial wavefunction coefficients}

The high-level sector helper should support both the physical binomial
distribution and the existing square-over-total-\(N_J\) convention.

\subsection*{Binomial distribution}

For the product state

\begin{equation}
\left(\frac{|\uparrow\rangle+|\downarrow\rangle}{\sqrt{2}}\right)^N,
\end{equation}

the probability of a group-resolved sector is

\begin{equation}
P(N_{J,1},N_{J,2})
=
\frac{
\binom{N_1}{N_{J,1}}
\binom{N_2}{N_{J,2}}
}{2^N}.
\end{equation}

Therefore the unnormalized amplitude should be

\begin{equation}
c_{N_{J,1},N_{J,2}}
\propto
\sqrt{
\binom{N_1}{N_{J,1}}
\binom{N_2}{N_{J,2}}
}.
\end{equation}

After restricting to the selected total-\(N_J\) window, normalize all amplitudes
so that

\begin{equation}
\sum_{N_{J,1},N_{J,2}}
|c_{N_{J,1},N_{J,2}}|^2
=1.
\end{equation}

\subsection*{Square-over-total-\(N_J\) distribution}

If \texttt{sector\_distribution="square"} is used, keep the old meaning: equal
total probability for each selected total \(N_J\) sector.

Do not assign equal probability to every group-resolved pair directly. Instead:

\begin{enumerate}
\item Give each selected total \(N_J\) equal total probability.
\item Split that probability among the valid \((N_{J,1},N_{J,2})\) pairs using
  the conditional binomial distribution.
\end{enumerate}

Thus

\begin{equation}
P(N_J)=\frac{1}{N_{\rm sectors}},
\end{equation}

where \(N_{\rm sectors}\) is the number of valid selected total-\(N_J\) sectors.

For fixed \(N_J\), use

\begin{equation}
P(N_{J,1},N_{J,2}\mid N_J)
=
\frac{
\binom{N_1}{N_{J,1}}
\binom{N_2}{N_{J,2}}
}{
\sum\limits_{\substack{a+b=N_J\\0\leq a\leq N_1\\0\leq b\leq N_2}}
\binom{N_1}{a}
\binom{N_2}{b}
}.
\end{equation}

Then

\begin{equation}
P(N_{J,1},N_{J,2})
=
P(N_J)
P(N_{J,1},N_{J,2}\mid N_J),
\end{equation}

and the sector amplitude is

\begin{equation}
c_{N_{J,1},N_{J,2}}
=
\sqrt{P(N_{J,1},N_{J,2})}.
\end{equation}

This preserves the old square distribution over total \(N_J\), while
distributing each total sector physically across the two coupling groups.

\subsection*{Flat group-resolved distribution}

Do not use a flat distribution over all generated \((N_{J,1},N_{J,2})\) pairs by
default. That changes the intended total-\(N_J\) distribution.

If needed later, add it as a separate explicit option such as

\begin{verbatim}
sector_distribution="square_group_resolved"
\end{verbatim}

but do not use it for the existing \texttt{"square"} behavior.

\section*{Group operators}

Build group operators using sparse tensor products.

For group 1,

\begin{equation}
J_{1,-}=J_-(N_{J,1})\otimes I_2,
\qquad
J_{1,+}=J_+(N_{J,1})\otimes I_2.
\end{equation}

For group 2,

\begin{equation}
J_{2,-}=I_1\otimes J_-(N_{J,2}),
\qquad
J_{2,+}=I_1\otimes J_+(N_{J,2}).
\end{equation}

The group \(x\)-operators are

\begin{equation}
J_{1,x}=\frac{J_{1,+}+J_{1,-}}{2},
\qquad
J_{2,x}=\frac{J_{2,+}+J_{2,-}}{2}.
\end{equation}

The group excited-state number operators are

\begin{equation}
N_{e,1}=N_e(N_{J,1})\otimes I_2,
\qquad
N_{e,2}=I_1\otimes N_e(N_{J,2}).
\end{equation}

The total excited-state number is

\begin{equation}
N_e=N_{e,1}+N_{e,2}.
\end{equation}

\section*{Inhomogeneous drive replacement}

Replace the homogeneous drive

\begin{equation}
\Omega J_x
\end{equation}

by

\begin{equation}
H_{\Omega}^{\rm inh}
=
\Omega(\omega_1J_{1,x}+\omega_2J_{2,x}).
\end{equation}

Equivalently,

\begin{equation}
H_{\Omega}^{\rm inh}
=
\frac{\Omega}{2}
\left[
\omega_1(J_{1,+}+J_{1,-})
+
\omega_2(J_{2,+}+J_{2,-})
\right].
\end{equation}

The detuning remains homogeneous:

\begin{equation}
H_\delta=-\delta(N_{e,1}+N_{e,2}).
\end{equation}

Thus, in the unshifted picture,

\begin{equation}
H^{\rm inh}
=
\Omega(\omega_1J_{1,x}+\omega_2J_{2,x})
-
\delta(N_{e,1}+N_{e,2}).
\end{equation}

\section*{Collective jump operator convention}

If only the drive is made inhomogeneous, the collective decay remains the same
total collective channel:

\begin{equation}
L=J_{1,-}+J_{2,-}.
\end{equation}

If the inhomogeneous coupling is also meant to modify the cavity/emission
coupling, use the weighted collective channel

\begin{equation}
A=\omega_1J_{1,-}+\omega_2J_{2,-}.
\end{equation}

This is still collective decay because there is one shared jump channel. It is
not two independent jump channels.

Do not replace this by

\begin{equation}
l_1=\omega_1J_{1,-},
\qquad
l_2=\omega_2J_{2,-},
\end{equation}

unless independent group-resolved decay is explicitly requested.

The dissipator is not linear in the jump operator:

\begin{equation}
\mathcal{D}[\omega_1J_{1,-}+\omega_2J_{2,-}]
\neq
\omega_1^2\mathcal{D}[J_{1,-}]
+
\omega_2^2\mathcal{D}[J_{2,-}].
\end{equation}

The cross terms encode collective emission into a shared environment.

\section*{Shifted-jump picture}

For the weighted collective model, define

\begin{equation}
A=\omega_1J_{1,-}+\omega_2J_{2,-}.
\end{equation}

The shifted jump operator is

\begin{equation}
l=A+i\frac{\Omega}{\Gamma}.
\end{equation}

Expanding the dissipator gives

\begin{equation}
\Gamma\mathcal{D}\left[A+i\frac{\Omega}{\Gamma}\right]\rho
=
\Gamma\mathcal{D}[A]\rho
-
i\left[\Omega\frac{A^\dagger+A}{2},\rho\right].
\end{equation}

Since

\begin{equation}
\frac{A^\dagger+A}{2}
=
\omega_1J_{1,x}+\omega_2J_{2,x},
\end{equation}

the shifted jump operator absorbs the full inhomogeneous drive Hamiltonian for
the weighted collective model.

Therefore,

\begin{equation}
H=
\Omega(\omega_1J_{1,x}+\omega_2J_{2,x})
-
\delta(N_{e,1}+N_{e,2}),
\qquad
l=A
\end{equation}

is equivalent to

\begin{equation}
H=-\delta(N_{e,1}+N_{e,2}),
\qquad
l=A+i\frac{\Omega}{\Gamma}.
\end{equation}

\section*{Coding implications}

\begin{itemize}
\item Add a high-level helper that expands \((N,dN,N_1,N_2,\omega_1)\) into all
  valid group-resolved sectors \((N_{J,1},N_{J,2})\).
\item Keep the low-level representation
  \(\{(\texttt{Nj1},\texttt{Nj2}):\texttt{coeff}\}\) internally.
\item Choose \(\omega_2=(N-N_1\omega_1)/N_2\) once from the physical group sizes,
  so \(N_1\omega_1+N_2\omega_2=N\).
\item Use the physical count condition \(N_{J,1}+N_{J,2}=N_J\) when mapping a
  target homogeneous-sector label \(N_J\) onto inhomogeneous sector pairs.
\item For \texttt{sector\_distribution="square"}, keep equal total probability
  over selected total \(N_J\) sectors and split each total sector using the
  conditional binomial distribution.
\item For \texttt{sector\_distribution="binomial"}, use amplitudes proportional
  to
  \(\sqrt{\binom{N_1}{N_{J,1}}\binom{N_2}{N_{J,2}}}\), normalized over the
  selected window.
\item Build separate group operators for each fixed pair \((N_{J,1},N_{J,2})\).
\item The Hilbert basis should track both group excitation numbers, for example
  \((n_{e,1},n_{e,2})\).
\item Use \(J_{1,\pm}\) only on group 1 and \(J_{2,\pm}\) only on group 2.
\item Use the same detuning \(\delta\) for both groups through
  \(-\delta(N_{e,1}+N_{e,2})\).
\item Use one collective jump operator, not two independent jumps.
\item Precompute sector operators, jump operators, effective generators, and
  propagators as in the homogeneous implementation.
\item Avoid rebuilding sparse Kronecker products or \(l^\dagger l\) inside the
  trajectory time loop.
\item Recover the homogeneous implementation when one active group is empty,
  e.g. \(N_{J,1}=0\) or \(N_{J,2}=0\), within numerical tolerance.
\end{itemize}

\section*{Mean-field steady-state residual check}

Add a standalone diagnostic for checking the inhomogeneous mean-field
steady-state equation from the paper.

Use the same fixed coupling weights as the simulation:

\begin{equation}
\omega_2=\frac{N-N_1\omega_1}{N_2}.
\end{equation}

Do not recompute \(\omega_2\) from instantaneous or sector-averaged values of
\(N_{J,1}\) and \(N_{J,2}\).

The steady-state phase equation is

\begin{equation}
\begin{aligned}
\frac{\Omega\omega_a}{2}
e^{-i\phi_a}
\sin\theta_a
=\;&
\frac{\delta}{2}
\sin\theta_a
\tan\theta_a
\\
&-
i
\frac{\Gamma\omega_a}{4}
e^{-i\phi_a}
\sin\theta_a
\left[
\omega_1N_{J,1}e^{i\phi_1}\sin\theta_1
+
\omega_2N_{J,2}e^{i\phi_2}\sin\theta_2
\right],
\qquad
a=1,2.
\end{aligned}
\end{equation}

This is the same equation as in the PDF, but written with \(\theta_a\) instead
of enforcing a common \(\theta_J\). Do not replace \(\theta_a\) by a common angle
in this diagnostic.

Define the left-hand side

\begin{equation}
{\rm LHS}_a
=
\frac{\Omega\omega_a}{2}
e^{-i\phi_a}
\sin\theta_a.
\end{equation}

Define the right-hand side

\begin{equation}
\begin{aligned}
{\rm RHS}_a
=\;&
\frac{\delta}{2}
\sin\theta_a
\tan\theta_a
\\
&-
i
\frac{\Gamma\omega_a}{4}
e^{-i\phi_a}
\sin\theta_a
\left[
\omega_1N_{J,1}e^{i\phi_1}\sin\theta_1
+
\omega_2N_{J,2}e^{i\phi_2}\sin\theta_2
\right].
\end{aligned}
\end{equation}

The residual should be

\begin{equation}
R_a
=
{\rm LHS}_a
-
{\rm RHS}_a.
\end{equation}

Explicitly,

\begin{equation}
\begin{aligned}
R_a
=\;&
\frac{\Omega\omega_a}{2}
e^{-i\phi_a}
\sin\theta_a
-
\frac{\delta}{2}
\sin\theta_a
\tan\theta_a
\\
&+
i
\frac{\Gamma\omega_a}{4}
e^{-i\phi_a}
\sin\theta_a
\left[
\omega_1N_{J,1}e^{i\phi_1}\sin\theta_1
+
\omega_2N_{J,2}e^{i\phi_2}\sin\theta_2
\right].
\end{aligned}
\end{equation}

For a steady state,

\begin{equation}
R_a\approx0,
\qquad
a=1,2.
\end{equation}

Here:

\begin{itemize}
\item \(a=1,2\) labels the coupling group.
\item \(N_{J,1}\) is the number of active-manifold atoms belonging to group 1.
\item \(N_{J,2}\) is the number of active-manifold atoms belonging to group 2.
\item \(\omega_1\) and \(\omega_2\) are the drive weights.
\item \(\Omega\) is the drive amplitude in the current phase.
\item \(\delta\) is the detuning in the current phase.
\item \(\Gamma\) is the collective decay rate.
\item \(\theta_a,\phi_a\) are the group-resolved active-manifold Bloch angles.
\end{itemize}

The residual plotting function should evaluate \(R_1(t)\) and \(R_2(t)\) from
snapshot data using the phase-local values of \(\Omega(t)\) and \(\delta(t)\),
then plot the absolute values of the real and imaginary parts of the two complex
residuals in a \(2\times2\) grid:

\begin{verbatim}
|Re R_1|    |Im R_1|
|Re R_2|    |Im R_2|
\end{verbatim}

\subsection*{Standalone residual plotting function}

Add a standalone post-processing function, without changing the core trajectory
implementation unless additional data must be stored.

Suggested function name:

\begin{verbatim}
plot_inhomogeneous_mfe_residuals(result, *, omega1, N1, N2,
                                 axes=None, output_path=None)
\end{verbatim}

The function should:

\begin{itemize}
\item read saved snapshot data from the trajectory or ensemble result;
\item determine the current phase for each snapshot;
\item use the phase-local values of \(\Omega(t)\) and \(\delta(t)\);
\item use the global \(\Gamma\) stored in the result;
\item compute group-resolved observables for each snapshot;
\item compute \(\theta_1(t),\phi_1(t)\) from group 1;
\item compute \(\theta_2(t),\phi_2(t)\) from group 2;
\item compute \(N_{J,1}(t)\) and \(N_{J,2}(t)\), or use the fixed
  group-resolved sector labels when evaluating within a sector;
\item compute \(R_1(t)\) and \(R_2(t)\);
\item plot the absolute values of the real and imaginary parts of \(R_1(t)\) and
  \(R_2(t)\) in a \(2\times2\) grid.
\end{itemize}

Use the \(2\times2\) grid as

\begin{verbatim}
|Re R_1|    |Im R_1|
|Re R_2|    |Im R_2|
\end{verbatim}

The plotted curves should be close to zero when the trajectory is near the
mean-field steady state.

\subsection*{Group and average angle plotting function}

Add another standalone post-processing function for plotting group-resolved and
average angles.

Suggested function name:

\begin{verbatim}
plot_inhomogeneous_group_angles(result, *, axes=None, output_path=None)
\end{verbatim}

This function should create a \(1\times2\) grid.

First panel:

\begin{itemize}
\item plot \(\theta_1(t)\);
\item plot \(\theta_2(t)\);
\item plot the correctly computed average angle \(\theta_{\rm avg}(t)\).
\end{itemize}

Second panel:

\begin{itemize}
\item plot \(\phi_1(t)\);
\item plot \(\phi_2(t)\);
\item plot the correctly computed average phase \(\phi_{\rm avg}(t)\).
\end{itemize}

The average-angle convention must follow
\texttt{docs/instructions/bloch\_vector\_averaging.tex}. In particular, do not
compute arithmetic averages such as \((\theta_1+\theta_2)/2\). Compute the
relevant group or sector moments first, normalize by \(N_{\rm active}\), and
only then convert the normalized Bloch direction to \(\theta\) and \(\phi\).

\subsection*{Data requirements}

The diagnostic functions require group-resolved observables. If these are not
already stored or computable from snapshots, add only the minimum needed support
to compute

\begin{equation}
J_{x,1},J_{y,1},J_{z,1},N_{e,1},
\qquad
J_{x,2},J_{y,2},J_{z,2},N_{e,2}.
\end{equation}

Do not modify unrelated plotting or trajectory code for this feature. The
residual check should be implemented as post-processing of saved snapshots
whenever possible.

\end{document}
